% ===========================================================
% results.tex
% ===========================================================
\section{Results and Validation}
\label{sec:results}

This section consolidates the quantitative validation of the
pipeline against two independent references: the published
$\lambda^{\mathrm{CEN}}$ feed and the PLEXOS solution
\texttt{.accdb} oracle.  The headline operating day is
2026-04-22; multi-day trends are summarised at the end.

\subsection{Aggregate accuracy}
Across the \num{2494} cells of 2026-04-22, the
\textit{uncorrected} pipeline output (using the empirical
$\lambda_b/\lambda_{\mathrm{ref}}$ loss factor) yields a mean
absolute error of $\$0.152$/MWh, a P90 of $\$0.471$/MWh, and a
maximum of $\$1.338$/MWh.  Filtering to non-wedge-flagged cells
(\num{2095} of \num{2494}) tightens these to $\$0.124$,
$\$0.417$, and $\$1.338$/MWh respectively.  The
\textit{PID-corrected} estimator that uses
$\mathrm{LF}^{\mathrm{PID}}_{b,t}$ from the PLEXOS \texttt{.accdb}
yields a mean of $\$0.294$/MWh and a max of $\$4.31$/MWh — slightly
worse on aggregate than the empirical estimator, because the
empirical estimator is mathematically self-consistent at $\lambda$
by construction.  The PID estimator is nonetheless preferred when
the use case requires a counterfactual interpretation
(e.g.\ what would $\hat\lambda$ have been at a candidate new
bus?), since the empirical estimator collapses on any new bus
without published $\lambda^{\mathrm{CEN}}$.

\subsection{Regime distribution}
Across the \num{2494} cells the regime distribution is
\texttt{renewable\_curtailment} \num{1073} (\SI{43.0}{\percent}),
\texttt{matched\_interior} \num{878} (\SI{35.2}{\percent}),
\texttt{matched\_capped\_pmax} \num{449} (\SI{18.0}{\percent}),
\texttt{matched\_above} \num{71} (\SI{2.8}{\percent}), and
\texttt{matched\_near\_pmin} \num{23} (\SI{0.9}{\percent}).  No
cells fall in the \texttt{no\_candidates}, \texttt{no\_cv\_match},
or \texttt{demand\_fail} regimes for this day.  The dominance of
\texttt{renewable\_curtailment} reflects the high mid-day solar
penetration of the SEN: across the four hours
\num{11}{:}\num{00}{-}\num{15}{:}\num{00} the
$\lambda^{\mathrm{CEN}}$ falls below $\$0.5$/MWh on most non-coastal
buses.

\subsection{Marginal-unit coherence}
By construction every bus in the same island shares the same
marginal-unit assignment.  A regression test reports
\num{0}/\num{474} islands where two buses disagree on the
marginal-unit name; this is the maximum-coherence outcome and
validates the clustering+picker composition.

\subsection{Per-bus MLF cross-check}
The PLEXOS-published per-bus MLFs read from the
2026-04-22 PID \texttt{.accdb} cover \num{247} buses
($\mathrm{LF}\in[0.92,1.05]$, mean $\approx0.99$, dominated by the
\SI{500}{\kV} centre-Santiago corridor).  After bridging by
\texttt{id\_info}, \num{26} of the \num{29} pipeline buses receive
a per-period MLF; the remaining \num{3} buses fall back to the
empirical estimator.  The mean cross-bus MLF spread on this day
is \num{0.072} (max $-$ min), in line with the typical $\pm$\SI{4}{\percent}
loss penalty observed in the SEN backbone.

\subsection{Pipeline runtime}
Table~\ref{tab:runtime} reports the runtime breakdown on a
warm cache (all CEN/SCVIC parquet present on disk) and on a cold
cache (no parquet, full network fetch).  The warm-cache run
completes in under \SI{6}{\second} dominated by the parallel
discovery of bar\_transfs and the per-island marginal-unit
identification loop.  The cold-cache run is dominated by
network rate-limiting on the per-bus cmg-real fetch, with
parallel \num{8}-thread fetch achieving an effective throughput of
roughly \num{0.4}~bus per second; one cold day for the full
\num{29}-bus catalogue requires approximately \SI{75}{\second}.

\begin{table}[t]
  \centering
  \caption{Pipeline runtime breakdown for 2026-04-22, on a
  warm cache (all parquet on disk) and on a cold cache (no parquet,
  full network fetch).}
  \label{tab:runtime}
  % Pipeline runtime breakdown (Section 6.5)
\begin{tabular}{l r r}
  \toprule
  Stage                                       & Warm cache [s] & Cold cache [s] \\
  \midrule
  Bus catalogue + bar\_transf discovery       & 0.10           &  3.5 \\
  Per-bus parallel \texttt{cmg-real} fetch (29 buses, 8 threads) & 0.20 & 64.0 \\
  \texttt{generaci\'on-real}                  & 0.05           &  1.8 \\
  \texttt{instrucciones-operacionales-cmg}    & 0.04           &  1.6 \\
  \texttt{unidades-generadoras}               & 0.05           &  1.4 \\
  \texttt{cv\_rpt} SCVIC parquet              & 0.03           &  0.8 \\
  PID \texttt{.accdb} MLF extraction (\texttt{mdbtools}) & 0.50 & 0.5 \\
  Alias index + per-id metadata               & 0.30           &  0.3 \\
  Island clustering ($96\!\times\!29$ cells)  & 0.40           &  0.4 \\
  Per-island marginal-unit picker (\num{474} keys) & 1.10      &  1.1 \\
  Output assembly (vectorised merge, \num{2494} rows) & 0.10   &  0.1 \\
  Loss-factor + $\varepsilon$ corrections     & 0.05           &  0.1 \\
  Parquet write                               & 0.05           &  0.1 \\
  \midrule
  \textbf{Total}                              & \textbf{$\approx$\,3.0–6.0\,s} & \textbf{$\approx$\,75\,s} \\
  \bottomrule
\end{tabular}
\\[3pt]
{\footnotesize Wall-clock measurements on a Linux x86\_64 laptop
(Ryzen 7~5800X, \SI{32}{GiB} RAM), single CPU thread except where
noted.  The cold-cache time is dominated by network rate-limiting
on the per-bar\_transf \texttt{cmg-real} fetch; with the
\texttt{cmg-online} bulk-paginated endpoint instead, the cold
total drops to roughly \SI{45}{\second} for the same \num{29} buses.}

\end{table}

\subsection{Cross-validation against the PLEXOS oracle}
Table~\ref{tab:oracle} summarises the agreement between the
pipeline-published quantities and the PLEXOS \texttt{.accdb} oracle
on 2026-04-22.  Three quantities are compared: the per-bus
$\lambda$ (the same quantity is published by both CEN SIP and the
PLEXOS solver, providing a publication-chain integrity check), the
per-unit dispatch (compared between SIP \texttt{generaci\'on-real}
and the PLEXOS per-unit dispatch), and the per-bus MLF (the only
quantity not in the SIP feed).  All three agree to within
publication noise, confirming that the methodology's reliance on
SIP feeds does not introduce a systematic bias relative to the
solver-side values.

\begin{table}[h]
  \centering
  \caption{Cross-validation against the PLEXOS PCP/PID
  \texttt{.accdb} oracle for 2026-04-22.  All three quantities
  agree to within publication noise.}
  \label{tab:oracle}
  \begin{tabular}{lrrr}
    \toprule
    Quantity & mean $|\Delta|$ & P90 & max \\
    \midrule
    $\lambda_b$ (USD/MWh) & 0.04 & 0.18 & 0.92 \\
    $p_{g,h}$ (MW)        & 0.6  & 2.7  & 14.3 \\
    $\mathrm{LF}_b$ (\%)  & 0.21 & 0.78 & 1.94 \\
    \bottomrule
  \end{tabular}
\end{table}
